\chapter{التنظيم العام للوحدة المركزية للحاسوب}

\begin{khulasa}[: ماذا سنتعلم في هذا الفصل؟]
\begin{itemize}
    \item ما الحاسوب، وكيف تطوّر عبر التاريخ.
    \item المكونات الأساسية: المعالج، الذاكرة، وحدات الإدخال والإخراج.
    \item نموذجا \textenglish{Von Neumann} و\textenglish{Harvard}، والفرق بينهما.
    \item أنواع الذاكرة، وكيف نحسب سعتها.
    \item الحافلات الثلاث: حافلة المعطيات، حافلة العناوين، حافلة التحكم.
\end{itemize}
\end{khulasa}

%================================================
\section{ما الحاسوب؟}
%================================================

\begin{tarif}
    \mustalah{الحاسوب}{Ordinateur}{Computer}: آلة إلكترونية تستقبل \textbf{معطيات}،
    وتعالجها حسب \textbf{برنامج} (قائمة من التعليمات)، ثم تعطي \textbf{نتائج}، ويمكنها أن تحفظ
    المعطيات والنتائج.
\end{tarif}

كل ما يفعله الحاسوب يمكن تلخيصه في أربعة أفعال: \textbf{يستقبل}، \textbf{يعالج}،
\textbf{يحفظ}، \textbf{يُخرج}.

لنأخذ مثالًا بسيطًا: تستعمل تطبيق الآلة الحاسبة في حاسوبك لتحسب $25 \times 4$.
\begin{enumerate}
    \item \textbf{الإدخال}: تكتب \textenglish{25}، ثم \textenglish{×}، ثم \textenglish{4} بلوحة المفاتيح.
          كل ضغطة تتحول إلى رمز رقمي يصل إلى الحاسوب.
    \item \textbf{الحفظ المؤقت}: يحفظ الحاسوب العددين في الذاكرة، لأنه يحتاج إليهما بعد قليل.
    \item \textbf{المعالجة}: يقرأ المعالج العددين، وينفذ تعليمة الضرب، فيحصل على \textenglish{100}.
    \item \textbf{الإخراج}: يرسل الحاسوب النتيجة إلى الشاشة فتراها.
\end{enumerate}
كل برنامج تستعمله، مهما كان معقدًا (لعبة، متصفح، تطبيق صور)، هو في الحقيقة ملايين من هذه
الخطوات الصغيرة، تُنفَّذ بسرعة كبيرة جدًا.

\begin{center}
\begin{tikzpicture}[node distance=1.4cm]
    \node[unit=brandGreen, minimum width=2.8cm, minimum height=1.2cm] (in) {الإدخال\\{\small لوحة المفاتيح، الفأرة}};
    \node[unit, minimum width=2.8cm, minimum height=1.2cm, left=of in] (cpu) {المعالجة\\{\small المعالج}};
    \node[unit=brandOrange, minimum width=2.8cm, minimum height=1.2cm, left=of cpu] (out) {الإخراج\\{\small الشاشة، الطابعة}};
    \node[unit=brandPurple, minimum width=2.8cm, minimum height=1.1cm, below=1cm of cpu] (mem) {الحفظ\\{\small الذاكرة، القرص}};
    \draw[sahmT] (in) -- (cpu);
    \draw[sahmT] (cpu) -- (out);
    \draw[sahmT, <->] (cpu) -- (mem);
\end{tikzpicture}
\end{center}

\begin{tashbih}
    الحاسوب يشبه \textbf{مطبخ مطعم}: المكونات (الخضر، اللحم) هي \textbf{المعطيات}، والوصفة هي
    \textbf{البرنامج}، والطباخ هو \textbf{المعالج}، والثلاجة هي \textbf{الذاكرة}، والطبق الجاهز
    هو \textbf{النتيجة}. الطباخ لا يقرر وحده: هو ينفذ الوصفة خطوة بخطوة.
\end{tashbih}

\subsection{العتاد والبرمجيات}
\begin{itemize}
    \item \mustalah{العتاد}{Matériel}{Hardware}: كل ما يمكن لمسه: المعالج، الذاكرة، الشاشة، الأسلاك.
    \item \mustalah{البرمجيات}{Logiciel}{Software}: البرامج التي تقول للعتاد ماذا يفعل:
          نظام التشغيل، المتصفح، الألعاب.
\end{itemize}
في هذا المقياس ندرس \textbf{العتاد من الداخل}، ونرى كيف ينفذ البرامج.

\subsection{طبقات الحاسوب}
يمكن أن نتخيل الحاسوب طبقات بعضها فوق بعض. كل طبقة تستعمل خدمات الطبقة التي تحتها، ولا تحتاج
أن تعرف تفاصيلها:

\begin{center}
\begin{tikzpicture}[l/.style={draw=#1, fill=#1!15, very thick, minimum width=9cm, minimum height=0.75cm, font=\small}]
    \node[l=brandGreen] at (0,1.8) {التطبيقات: المتصفح، الألعاب، معالج النصوص};
    \node[l=brandPurple] at (0,0.9) {نظام التشغيل: \textenglish{Windows}، \textenglish{Linux}، \textenglish{Android}};
    \node[l=brandOrange] at (0,0) {لغة الآلة ولغة التجميع: التعليمات التي يفهمها المعالج};
    \node[l=brandTeal] at (0,-0.9) {العتاد: المعالج، الذاكرة، الحافلات، الدارات الإلكترونية};
    \draw[sahmT] (5,-1.1) -- (5,2) node[midway, right, font=\small] {من الآلة إلى الإنسان};
\end{tikzpicture}
\end{center}

في هذا الدرس نشتغل على \textbf{الطبقتين السفليين}: العتاد، ولغة التجميع التي تتكلم معه مباشرة.

\subsection{كيف يعمل برنامج؟}
عندما تضغط مرتين على أيقونة برنامج، تحدث هذه المراحل:
\begin{enumerate}
    \item البرنامج محفوظ على \textbf{القرص} (\textenglish{SSD} أو \textenglish{HDD}) في شكل ملف من الأصفار والآحاد.
    \item نظام التشغيل \textbf{ينسخه في الذاكرة المركزية} \textenglish{(RAM)}، لأن القرص بطيء جدًا والمعالج لا يقرأ منه مباشرة.
    \item المعالج يقرأ التعليمات من الذاكرة \textbf{واحدة بعد الأخرى} وينفذها.
    \item عندما تغلق البرنامج، تُحرَّر مساحته في الذاكرة لبرامج أخرى.
\end{enumerate}

%================================================
\section{تاريخ الحواسيب باختصار}
%================================================

لم يظهر الحاسوب فجأة. هذه أهم المحطات:

\begin{center}
\begin{tikzpicture}[
    yr/.style={circle, fill=brandTeal, minimum size=0.35cm, inner sep=0pt},
    txt/.style={text width=2.7cm, align=center, font=\small}]
    \draw[brandTeal, line width=3pt] (0,0) -- (-14,0);
    \foreach \x/\y/\t/\p in {
        0/1642/{آلة \textenglish{Pascal}\\ للجمع والطرح}/above,
        -2/1837/{آلة \textenglish{Babbage}\\ فكرة البرنامج}/below,
        -4/1945/{\textenglish{ENIAC}\\ مصابيح مفرغة}/above,
        -6/1947/{اختراع\\ \textbf{الترانزستور}}/below,
        -8/1958/{\textbf{الدارة المندمجة}\\ \textenglish{(IC)}}/above,
        -10/1971/{أول معالج دقيق\\ \textenglish{Intel 4004}}/below,
        -12/1978/{المعالج\\ \textenglish{Intel 8086}}/above,
        -14/2005/{المعالجات\\ \textbf{متعددة الأنوية}}/below} {
        \node[yr] (n) at (\x, 0) {};
        \node[txt, \p=0.25cm of n] {\textbf{\textenglish{\y}}\\ \t};
    }
\end{tikzpicture}
\end{center}

\begin{table}[h]
\centering
\begin{tabular}{|c|c|c|c|}
    \hline
    \ras{الجيل} & \ras{الفترة} & \ras{التقنية الأساسية} & \ras{مثال} \\ \hline
    الأول   & \textenglish{1945--1955} & المصابيح المفرغة & \textenglish{ENIAC} (\textenglish{30} طنًا!) \\ \hline
    الثاني  & \textenglish{1955--1965} & الترانزستور & \textenglish{IBM 7090} \\ \hline
    الثالث  & \textenglish{1965--1980} & الدارات المندمجة & \textenglish{IBM 360} \\ \hline
    الرابع  & \textenglish{1980--} الآن & المعالج الدقيق (\textenglish{VLSI}) & الحاسوب الشخصي، الهاتف \\ \hline
\end{tabular}
\end{table}

كل جيل جديد جاء بتقنية تجعل المكونات \textbf{أصغر} و\textbf{أسرع} و\textbf{أقل استهلاكًا للطاقة}:
\begin{itemize}
    \item \textbf{المصابيح المفرغة} كانت كبيرة، تسخن كثيرًا، وتحترق باستمرار؛ كان يجب تبديل بعضها كل يوم تقريبًا.
    \item \textbf{الترانزستور} يقوم بالعمل نفسه (مفتاح كهربائي يفتح ويغلق)، لكنه أصغر بمئات المرات ولا يحترق.
    \item \textbf{الدارة المندمجة} تضع آلاف الترانزستورات على قطعة صغيرة واحدة من السيليسيوم، بدل توصيلها بالأسلاك.
    \item \textbf{المعالج الدقيق} يضع المعالج \textbf{كله} في شريحة واحدة. من هنا ولد الحاسوب الشخصي ثم الهاتف الذكي.
\end{itemize}

\begin{tarif}
    \mustalah{قانون مور}{Loi de Moore}{Moore's law}: لاحظ \textenglish{Gordon Moore} (أحد مؤسسي \textenglish{Intel})
    سنة 1965 أن عدد الترانزستورات في الشريحة \textbf{يتضاعف تقريبًا كل سنتين}. بقيت هذه الملاحظة صحيحة
    لأكثر من خمسين سنة: من \textenglish{2300} ترانزستور في \textenglish{4004} إلى عشرات المليارات اليوم.
\end{tarif}

\begin{hal}
    كان \textenglish{ENIAC} يملأ قاعة كاملة ويستهلك \textenglish{150} كيلوواط، بينما هاتفك
    أقوى منه بملايين المرات ويعمل ببطارية صغيرة. السبب: الترانزستورات أصبحت صغيرة جدًا،
    حتى صار المعالج الواحد يحتوي على مليارات منها.
\end{hal}

%================================================
\section{تذكير سريع: كيف يعدّ الحاسوب؟}
%================================================

الحاسوب لا يعرف إلا حالتين كهربائيتين: \textbf{يمر التيار} أو \textbf{لا يمر}. نكتبهما
\textenglish{1} و\textenglish{0}.

\begin{tarif}
    \begin{itemize}
        \item \mustalah{البت}{Bit}{Bit}: أصغر وحدة معلومة، قيمته \textenglish{0} أو \textenglish{1}.
        \item \mustalah{البايت}{Octet}{Byte}: مجموعة من \textbf{8 بتات}.
        \item \mustalah{الكلمة}{Mot}{Word}: عدد البتات التي يعالجها المعالج مرة واحدة
              (\textenglish{16} بتًا في المعالج \textenglish{8086}).
    \end{itemize}
\end{tarif}

نكتب الأعداد في ثلاثة أنظمة: العشري (الأساس 10)، الثنائي (الأساس 2)، والست عشري (الأساس 16،
أرقامه من \textenglish{0} إلى \textenglish{9} ثم \textenglish{A} إلى \textenglish{F}).

\begin{center}
\small
\begin{tabular}{|c|c|c||c|c|c|}
    \hline
    \ras{عشري} & \ras{ثنائي} & \ras{ست عشري} & \ras{عشري} & \ras{ثنائي} & \ras{ست عشري} \\ \hline
    \textenglish{0} & \textenglish{0000} & \textenglish{0} & \textenglish{8}  & \textenglish{1000} & \textenglish{8} \\ \hline
    \textenglish{1} & \textenglish{0001} & \textenglish{1} & \textenglish{9}  & \textenglish{1001} & \textenglish{9} \\ \hline
    \textenglish{2} & \textenglish{0010} & \textenglish{2} & \textenglish{10} & \textenglish{1010} & \textenglish{A} \\ \hline
    \textenglish{3} & \textenglish{0011} & \textenglish{3} & \textenglish{11} & \textenglish{1011} & \textenglish{B} \\ \hline
    \textenglish{4} & \textenglish{0100} & \textenglish{4} & \textenglish{12} & \textenglish{1100} & \textenglish{C} \\ \hline
    \textenglish{5} & \textenglish{0101} & \textenglish{5} & \textenglish{13} & \textenglish{1101} & \textenglish{D} \\ \hline
    \textenglish{6} & \textenglish{0110} & \textenglish{6} & \textenglish{14} & \textenglish{1110} & \textenglish{E} \\ \hline
    \textenglish{7} & \textenglish{0111} & \textenglish{7} & \textenglish{15} & \textenglish{1111} & \textenglish{F} \\ \hline
\end{tabular}
\end{center}

\begin{khatwa}[: تحويل \textenglish{2F}\textsubscript{\textenglish{H}} إلى الثنائي وإلى العشري]
    \begin{enumerate}
        \item كل رقم ست عشري يساوي \textbf{4 بتات}: \textenglish{2 = 0010} و\textenglish{F = 1111}.
        \item نضعهما جنبًا إلى جنب: $\hex{2F} = \bin{0010\,1111}$.
        \item إلى العشري: $2 \times 16 + 15 = 32 + 15 = 47$.
    \end{enumerate}
\end{khatwa}

\begin{khatwa}[: تحويل العدد الثنائي \textenglish{1011} إلى العشري]
    كل بت له «وزن» هو قوة من قوى 2، تبدأ من اليمين بـ$2^0 = 1$:
    \begin{center}
    \begin{tabular}{|c|c|c|c|c|}
        \hline \ras{البت} & 1 & 0 & 1 & 1 \\ \hline
        \ras{الوزن} & $2^3 = 8$ & $2^2 = 4$ & $2^1 = 2$ & $2^0 = 1$ \\ \hline
        \ras{البت × الوزن} & 8 & 0 & 2 & 1 \\ \hline
    \end{tabular}
    \end{center}
    نجمع: $8 + 0 + 2 + 1 = 11$. إذن $\bin{1011} = 11$.
\end{khatwa}

\begin{khatwa}[: تحويل العدد العشري 13 إلى الثنائي (القسمة المتتالية على 2)]
    \begin{enumerate}
        \item $13 \div 2 = 6$، والباقي \textbf{1}.
        \item $6 \div 2 = 3$، والباقي \textbf{0}.
        \item $3 \div 2 = 1$، والباقي \textbf{1}.
        \item $1 \div 2 = 0$، والباقي \textbf{1}. نتوقف لأن الحاصل أصبح 0.
        \item نقرأ البواقي \textbf{من الأخير إلى الأول}: \bin{1101}. تحقق: $8 + 4 + 0 + 1 = 13$.
    \end{enumerate}
\end{khatwa}

\begin{tashbih}
    لماذا نستعمل الست عشري؟ لأن العدد الثنائي طويل جدًا ومتعب للعين: \bin{1010\,1111\,0011\,0110}.
    الست عشري يختصر كل 4 بتات في رقم واحد، فيصبح العدد \hex{AF36}. هو «الاختزال» الذي يستعمله المبرمجون
    لكتابة الثنائي بسرعة، كما نكتب «د.» بدل «دكتور».
\end{tashbih}

\begin{intabih}
    في لغة التجميع نكتب حرف \code{h} بعد العدد الست عشري (\code{2Fh})، وحرف \code{b} بعد العدد
    الثنائي (\code{00101111b}). العدد دون حرف هو عدد عشري.
\end{intabih}

\subsection{وحدات قياس الذاكرة}
\begin{center}
\begin{tabular}{|c|c|c|}
    \hline
    \ras{الوحدة} & \ras{الرمز} & \ras{القيمة} \\ \hline
    كيلوبايت  & \textenglish{KB} & $2^{10} = 1024$ بايت \\ \hline
    ميغابايت  & \textenglish{MB} & $2^{20} = 1\,048\,576$ بايت \\ \hline
    غيغابايت  & \textenglish{GB} & $2^{30}$ بايت (حوالي مليار) \\ \hline
    تيرابايت  & \textenglish{TB} & $2^{40}$ بايت \\ \hline
\end{tabular}
\end{center}

%================================================
\section{المكونات الأساسية للحاسوب}
%================================================

داخل كل حاسوب، من أكبر خادم إلى أصغر ساعة ذكية، نجد ثلاثة أجزاء أساسية:

\begin{center}
\begin{tikzpicture}
    \node[unit, minimum width=4cm, minimum height=2.6cm] (cpu) at (0,0) {};
    \node[font=\bfseries] at (0, 1) {المعالج \textenglish{(CPU)}};
    \node[unit=brandOrange, minimum width=1.6cm, font=\small] at (-0.95, 0) {\textenglish{UC}\\وحدة التحكم};
    \node[unit=brandGreen, minimum width=1.6cm, font=\small] at (0.95, 0) {\textenglish{UAL}\\الحساب};
    \node[draw=brandDark, fill=white, font=\small, minimum width=3.4cm] at (0,-0.9) {السجلات};
    \node[unit=brandPurple, minimum width=3.4cm, minimum height=2.6cm] (mem) at (-6,0) {الذاكرة المركزية\\{\small \textenglish{RAM} و\textenglish{ROM}}};
    \node[unit=brandGreen, minimum width=3.4cm, minimum height=2.6cm] (io) at (6,0) {وحدات\\الإدخال والإخراج};
    \draw[sahmT, <->] (cpu) -- (mem);
    \draw[sahmT, <->] (cpu) -- (io);
\end{tikzpicture}
\end{center}

\begin{enumerate}
    \item \mustalah{المعالج}{Processeur}{Processor, CPU}: «دماغ» الحاسوب. يقرأ التعليمات وينفذها.
    \item \mustalah{الذاكرة المركزية}{Mémoire centrale}{Main memory}: تحفظ البرنامج
          الذي يعمل الآن ومعطياته.
    \item \mustalah{وحدات الإدخال والإخراج}{Périphériques d'E/S}{I/O devices}: تسمح
          للحاسوب بالتواصل مع العالم الخارجي.
\end{enumerate}

\begin{khatwa}[: رحلة حرف من لوحة المفاتيح إلى الشاشة]
    تضغط على المفتاح \textenglish{A} في برنامج كتابة:
    \begin{enumerate}
        \item لوحة المفاتيح (\textbf{إدخال}) ترسل رمز المفتاح إلى الحاسوب.
        \item المعالج يُقاطَع (يتوقف لحظة عن عمله)، فيقرأ الرمز، ويحفظه في \textbf{الذاكرة} داخل النص الذي تكتبه.
        \item برنامج الكتابة يطلب من المعالج «رسم» الحرف: يحسب المعالج النقاط الضوئية التي تشكّل الحرف.
        \item هذه النقاط تُرسل إلى بطاقة الشاشة، ثم إلى الشاشة (\textbf{إخراج}).
    \end{enumerate}
    كل هذا يحدث في جزء صغير جدًا من الثانية، فتشعر أن الحرف ظهر فورًا.
\end{khatwa}

\begin{tarif}
    \mustalah{الوحدة المركزية}{Unité centrale}{Central unit}: المعالج مع الذاكرة المركزية.
    هي الجزء الذي ينفذ البرامج فعلًا.
\end{tarif}

%================================================
\section{النموذجان الأساسيان: فون نيومان وهارفارد}
%================================================

\subsection{نموذج فون نيومان \textenglish{(Von Neumann)}}

في سنة 1945 اقترح العالم \textenglish{John von Neumann} فكرة بسيطة وعبقرية: نحفظ
\textbf{البرنامج} و\textbf{المعطيات} معًا في \textbf{الذاكرة نفسها}. قبل ذلك، كان تغيير البرنامج
يتطلب إعادة توصيل الأسلاك يدويًا!

حسب فون نيومان، يتكون الحاسوب من \textbf{خمس وحدات}: وحدة الإدخال، وحدة الإخراج، الذاكرة، وحدة
الحساب والمنطق، ووحدة التحكم. والفكرة الأهم هي \mustalah{البرنامج المخزَّن}{Programme enregistré}{Stored program}:
التعليمات أعداد مثل المعطيات تمامًا، فنستطيع أن نحفظها في الذاكرة نفسها، ونقرأها ونغيّرها بالطريقة نفسها.

\begin{center}
\begin{tikzpicture}
    \node[unit, minimum width=4.2cm, minimum height=2.4cm] (cpu) at (0,0) {};
    \node[font=\bfseries] at (0,0.85) {المعالج};
    \node[unit=brandOrange, font=\small, minimum width=1.7cm] at (-1,-0.15) {وحدة التحكم};
    \node[unit=brandGreen, font=\small, minimum width=1.7cm] at (1,-0.15) {\textenglish{UAL}};
    \node[unit=brandPurple, minimum width=4.2cm, minimum height=2.4cm, align=center] (mem) at (-7,0) {};
    \node[font=\bfseries] at (-7,0.85) {ذاكرة واحدة};
    \node[draw=brandDark, fill=brandTeal!20, font=\small, minimum width=3.4cm] at (-7,0.1) {البرنامج (التعليمات)};
    \node[draw=brandDark, fill=brandOrange!20, font=\small, minimum width=3.4cm] at (-7,-0.6) {المعطيات};
    \draw[brandDark, line width=5pt, -{Stealth[length=4mm]}] (-4.8,0) -- (-2.2,0);
    \draw[brandDark, line width=5pt, -{Stealth[length=4mm]}] (-2.2,0) -- (-4.8,0);
    \node[font=\small, above] at (-3.5,0.2) {حافلة واحدة};
    \node[unit=brandGreen, minimum width=2.6cm, minimum height=1.2cm] (io) at (5,0) {الإدخال\\والإخراج};
    \draw[sahmT, <->] (cpu) -- (io);
\end{tikzpicture}
\end{center}

\begin{itemize}
    \item \textbf{ميزته}: بسيط ومرن. نغيّر البرنامج بمجرد تحميل برنامج آخر في الذاكرة.
    \item \textbf{عيبه}: حافلة واحدة للتعليمات والمعطيات، فلا يمكن قراءة تعليمة ومعطاة في
          اللحظة نفسها. هذا الاختناق يسمى \mustalah{عنق زجاجة فون نيومان}{Goulot de Von Neumann}{Von Neumann bottleneck}.
\end{itemize}

\begin{tashbih}
    تخيّل مكتبة لها \textbf{باب واحد ضيق}: الطلبة الذين يدخلون لإرجاع الكتب (المعطيات) والذين يدخلون
    لأخذ التعليمات (البرنامج) يمرون كلهم من الباب نفسه، واحدًا واحدًا. مهما كان الموظفون سريعين في
    الداخل، يبقى الباب هو ما يحدد السرعة. هذا هو عنق الزجاجة.
\end{tashbih}

\subsection{نموذج هارفارد \textenglish{(Harvard)}}

في هذا النموذج نستعمل \textbf{ذاكرتين منفصلتين}: واحدة للتعليمات وأخرى للمعطيات،
ولكل واحدة حافلتها الخاصة.

\begin{center}
\begin{tikzpicture}
    \node[unit, minimum width=3.4cm, minimum height=2.2cm] (cpu) at (0,0) {المعالج};
    \node[unit=brandTeal, minimum width=3.6cm, minimum height=1.6cm] (im) at (-6.5,0) {ذاكرة التعليمات\\{\small البرنامج}};
    \node[unit=brandOrange, minimum width=3.6cm, minimum height=1.6cm] (dm) at (6.5,0) {ذاكرة المعطيات\\{\small المتغيرات}};
    \draw[tealDeep, line width=4pt, -{Stealth[length=4mm]}] (im) -- (cpu) node[midway, above, font=\small, lawn=brandDark] {حافلة التعليمات};
    \draw[orangeDeep, line width=4pt, {Stealth[length=4mm]}-{Stealth[length=4mm]}] (cpu) -- (dm) node[midway, above, font=\small, lawn=brandDark] {حافلة المعطيات};
\end{tikzpicture}
\end{center}

\begin{itemize}
    \item \textbf{ميزته}: أسرع، لأن المعالج يقرأ التعليمة التالية ويقرأ المعطاة في الوقت نفسه.
    \item \textbf{عيبه}: أعقد وأغلى، والمساحة المخصصة لكل ذاكرة ثابتة.
\end{itemize}

\begin{tashbih}
    في نموذج هارفارد للمكتبة \textbf{بابان}: باب لأخذ التعليمات، وباب للمعطيات. الطابوران يتقدمان في
    الوقت نفسه، فيمر الناس أسرع. لكن بناء بابين أغلى، وإذا امتلأ أحد الجانبين لا نستطيع استعمال مساحة الجانب الآخر.
\end{tashbih}

\subsection{مقارنة واستعمالات}
\begin{center}
\begin{tabularx}{\linewidth}{|>{\centering\arraybackslash}p{3cm}|C|C|}
    \hline
    \ras{} & \ras{فون نيومان} & \ras{هارفارد} \\ \hline
    الذاكرة   & واحدة للبرنامج والمعطيات & ذاكرتان منفصلتان \\ \hline
    الحافلات  & حافلة واحدة مشتركة & حافلتان مستقلتان \\ \hline
    السرعة    & أبطأ (عنق الزجاجة) & أسرع (قراءتان معًا) \\ \hline
    التعقيد   & بسيط ورخيص & أعقد وأغلى \\ \hline
    الاستعمال & الحواسيب الشخصية، الخوادم & المتحكمات الدقيقة (\textenglish{Arduino})، معالجات الإشارة \textenglish{(DSP)} \\ \hline
\end{tabularx}
\end{center}

\begin{hal}
    المعالجات الحديثة في حاسوبك تجمع بين النموذجين: الذاكرة الرئيسية \textenglish{(RAM)} واحدة
    (فون نيومان)، لكن داخل المعالج توجد ذاكرتان صغيرتان سريعتان منفصلتان: واحدة للتعليمات
    وأخرى للمعطيات (هارفارد). نسمي ذلك \textbf{هارفارد المعدّل}، وسنعود إليه في الفصل الخامس.
\end{hal}

%================================================
\section{المعالج: الدور والمكونات}
%================================================

\begin{tarif}
    \mustalah{المعالج}{Processeur, UCT}{CPU}: الدارة التي تنفذ تعليمات البرنامج واحدة تلو الأخرى.
\end{tarif}

يتكون المعالج من ثلاثة أجزاء رئيسية (سندرسها بالتفصيل في الفصل الثاني):

\begin{center}
\begin{tabularx}{\linewidth}{|>{\centering\arraybackslash}p{4.2cm}|R|R|}
    \hline
    \ras{الجزء} & \ras{دوره} & \ras{التشبيه} \\ \hline
    وحدة التحكم \textenglish{(UC)} & تقرأ التعليمة، تفهمها، وتعطي الأوامر لباقي الأجزاء & قائد الأوركسترا \\ \hline
    وحدة الحساب والمنطق \textenglish{(UAL)} & تنفذ العمليات الحسابية (\textenglish{+ −}) والمنطقية (و، أو) & الآلة الحاسبة \\ \hline
    السجلات \textenglish{(Registres)} & ذاكرة صغيرة جدًا وسريعة جدًا داخل المعالج & ورقة المسودة على المكتب \\ \hline
\end{tabularx}
\end{center}

\subsection{الساعة وسرعة المعالج}
يعمل المعالج على إيقاع \mustalah{ساعة}{Horloge}{Clock} ترسل نبضات منتظمة. في كل نبضة يتقدم
المعالج خطوة صغيرة في عمله.

\begin{tarif}
    \textbf{التردد} $f$: عدد النبضات في الثانية، ويقاس بالهرتز \textenglish{(Hz)}.
    \textbf{الدورة} $T$: مدة نبضة واحدة، و$T = \dfrac{1}{f}$.
\end{tarif}

\begin{khatwa}[: معالج تردده \textenglish{3 GHz}]
    \begin{enumerate}
        \item $f = 3\ \text{GHz} = 3 \times 10^{9}$ نبضة في الثانية.
        \item $T = \dfrac{1}{3 \times 10^{9}} \approx 0.33 \times 10^{-9}$ ثانية، أي حوالي
              \textenglish{0.33} نانوثانية.
        \item في الوقت الذي ترمش فيه عينك (\textenglish{0.1} ثانية تقريبًا) يقوم هذا المعالج بـ
              \textenglish{300} مليون دورة!
    \end{enumerate}
\end{khatwa}

\begin{intabih}
    التردد وحده لا يحدد سرعة المعالج. ما يهمنا هو \textbf{عدد التعليمات المنفذة في الثانية}.
    معالج بـ\textenglish{3 GHz} يحتاج 4 دورات لكل تعليمة ينفذ $\frac{3 \times 10^9}{4} = 750$ مليون تعليمة
    في الثانية، بينما معالج بـ\textenglish{2 GHz} ينفذ تعليمتين في كل دورة ينفذ $4 \times 10^9$ تعليمة، فهو
    أسرع رغم أن تردده أقل.
\end{intabih}

%================================================
\section{الذاكرة}
%================================================

\subsection{كيف تبدو الذاكرة من الداخل؟}
نتخيل الذاكرة مثل \textbf{عمارة كبيرة من الصناديق}. لكل صندوق رقم يسمى \textbf{العنوان}،
وفي كل صندوق نحفظ \textbf{بايتًا واحدًا}.

\begin{center}
\begin{tikzpicture}[cellx/.style={draw=brandDark, fill=white, minimum width=2.4cm, minimum height=0.6cm, font=\codefont},
    addr/.style={font=\codefont\small, lawn=tealDeep}]
    \foreach \a/\v [count=\i from 0] in {0000/3A, 0001/FF, 0002/07, 0003/41, 0004/00} {
        \node[cellx] (c\i) at (0, -0.6*\i) {\textenglish{\v}};
        \node[addr, right=0.2cm of c\i] {\textenglish{\a}};
    }
    \node[font=\small\bfseries] at (2.15, 0.6) {العنوان};
    \node[font=\small\bfseries] at (0, 0.6) {المحتوى};
    \node[font=\small, text width=6.5cm, align=right] at (-5.5,-1.2)
        {العنوان \textenglish{0003} يحتوي على \textenglish{41}\textsubscript{\textenglish{H}}.\\
         العنوان لا يتغير أبدًا، أما المحتوى فيمكن تغييره.};
\end{tikzpicture}
\end{center}

\begin{tashbih}
    العنوان يشبه \textbf{رقم المنزل} في شارع، والمحتوى يشبه \textbf{الشخص الذي يسكن فيه}.
    رقم المنزل ثابت، أما السكان فيمكن أن يتغيروا.
\end{tashbih}

\subsection{عمليتان فقط: القراءة والكتابة}
\begin{itemize}
    \item \textbf{القراءة}: يرسل المعالج عنوانًا، فتعيد الذاكرة المحتوى (المحتوى لا يتغير).
    \item \textbf{الكتابة}: يرسل المعالج عنوانًا ومعطاة، فتحفظ الذاكرة المعطاة مكان القديمة.
\end{itemize}

\begin{mithal}
    الخانة \textenglish{0002} تحتوي على \hex{07}. ينفذ المعالج «اكتب \hex{55} في \textenglish{0002}».
    بعدها تحتوي الخانة على \hex{55}، و\hex{07} ضاعت نهائيًا. لهذا يجب أن ننتبه أين نكتب: الكتابة
    في المكان الخطأ تمحو معطيات مهمة.
\end{mithal}

\begin{tarif}
    \mustalah{ذاكرة الوصول العشوائي}{Mémoire à accès aléatoire}{Random Access Memory}: «عشوائي» هنا تعني
    أننا نستطيع الوصول إلى \textbf{أي خانة مباشرة} وفي الزمن نفسه، دون المرور بالخانات التي قبلها، عكس
    شريط الكاسيت القديم الذي يجب لفّه حتى نصل إلى الأغنية.
\end{tarif}

\subsection{أنواع الذاكرة}
\begin{center}
\begin{tikzpicture}[level distance=1.4cm,
    level 1/.style={sibling distance=7cm},
    level 2/.style={sibling distance=2.3cm},
    every node/.append style={unit, font=\small, minimum width=1.9cm},
    edge from parent/.style={draw, brandMuted, thick}]
    \node[unit=brandDark, fill=brandDark!10] {الذاكرة المركزية}
        child {node[unit=brandOrange] {\textenglish{RAM}\\{\scriptsize قراءة وكتابة}}
            child {node[unit=brandOrange] {\textenglish{SRAM}\\{\scriptsize سريعة}}}
            child {node[unit=brandOrange] {\textenglish{DRAM}\\{\scriptsize رخيصة}}}}
        child {node[unit=brandTeal] {\textenglish{ROM}\\{\scriptsize قراءة فقط}}
            child {node {\textenglish{PROM}}}
            child {node {\textenglish{EPROM}}}
            child {node {\textenglish{EEPROM}\\{\scriptsize و\textenglish{Flash}}}}};
\end{tikzpicture}
\end{center}

\begin{center}
\small
\begin{tabularx}{\linewidth}{|>{\centering\arraybackslash}p{2.2cm}|R|>{\centering\arraybackslash}p{2.4cm}|R|}
    \hline
    \ras{النوع} & \ras{الوصف} & \ras{عند انقطاع التيار} & \ras{الاستعمال} \\ \hline
    \textenglish{SRAM} & ترانزستورات فقط، سريعة جدًا وغالية & تُمحى & الذاكرة المخبئية داخل المعالج \\ \hline
    \textenglish{DRAM} & مكثّف صغير لكل بت، يجب «تنشيطه» باستمرار & تُمحى & الذاكرة الرئيسية (شرائط \textenglish{RAM}) \\ \hline
    \textenglish{ROM}  & تُكتب مرة واحدة في المصنع & تبقى & برامج الإقلاع القديمة \\ \hline
    \textenglish{EPROM} & تُمحى بالأشعة فوق البنفسجية ثم تُعاد كتابتها & تبقى & النماذج الأولية \\ \hline
    \textenglish{EEPROM / Flash} & تُمحى كهربائيًا & تبقى & \textenglish{BIOS/UEFI}، المفاتيح \textenglish{USB}، \textenglish{SSD} \\ \hline
\end{tabularx}
\end{center}

\begin{tashbih}
    \textbf{\textenglish{SRAM}} مثل مفتاح الضوء: يبقى في وضعه (مشتعل أو مطفأ) ما دام التيار موجودًا.
    أما \textbf{\textenglish{DRAM}} فمثل كأس فيه ثقب صغير: الماء (الشحنة الكهربائية) يتسرب ببطء، فيجب إعادة ملئه
    آلاف المرات في الثانية. هذه العملية تسمى \textbf{التنشيط} \textenglish{(refresh)}. لكن الكأس أصغر بكثير من المفتاح،
    لذلك نستطيع وضع مليارات منها في شريط ذاكرة رخيص.
\end{tashbih}

\begin{hal}
    لماذا يحتاج الحاسوب إلى \textenglish{ROM} أو \textenglish{Flash}؟ عند تشغيله تكون \textenglish{RAM} فارغة تمامًا.
    يجب أن يجد المعالج برنامجًا أول يبدأ به: برنامج الإقلاع (\textenglish{BIOS} أو \textenglish{UEFI})، المحفوظ في
    ذاكرة لا تُمحى. هذا البرنامج يفحص العتاد، ثم يبحث عن نظام التشغيل على القرص ويحمّله في \textenglish{RAM}.
\end{hal}

\begin{tarif}
    ذاكرة \mustalah{متطايرة}{Volatile}{Volatile} تفقد محتواها عند انقطاع التيار (\textenglish{RAM})،
    وذاكرة \mustalah{غير متطايرة}{Non volatile}{Non-volatile} تحتفظ به (\textenglish{ROM}، \textenglish{Flash}).
\end{tarif}

\subsection{حساب سعة الذاكرة}
إذا كان للمعالج $n$ خطًا (سلكًا) للعنوان، فيمكنه أن يكتب $2^n$ عنوانًا مختلفًا.

\begin{tarif}
    $\text{عدد الخانات} = 2^{n}$ \qquad و\qquad
    $\text{السعة} = 2^{n} \times \text{حجم الخانة}$
\end{tarif}

\begin{khatwa}[: سعة ذاكرة المعالج \textenglish{8086}]
    للمعالج \textenglish{8086} حافلة عناوين من \textbf{20 بتًا}، وكل خانة تحتوي على بايت واحد.
    \begin{enumerate}
        \item عدد العناوين: $2^{20}$.
        \item $2^{20} = 2^{10} \times 2^{10} = 1024 \times 1024 = 1\,048\,576$ عنوانًا.
        \item السعة $= 2^{20} \times 1$ بايت $= 1$ ميغابايت \textenglish{(1 MB)}.
    \end{enumerate}
\end{khatwa}

\begin{khatwa}[: كم خطًا للعنوان نحتاج لذاكرة \textenglish{64 KB}؟]
    \begin{enumerate}
        \item $64\ \text{KB} = 64 \times 1024 = 2^{6} \times 2^{10} = 2^{16}$ بايت.
        \item إذن $2^{n} = 2^{16}$، أي $n = 16$ خطًا.
    \end{enumerate}
\end{khatwa}

%================================================
\section{الحافلات}
%================================================

\begin{tarif}
    \mustalah{الحافلة}{Bus}{Bus}: مجموعة من الأسلاك المتوازية، تنقل المعلومات بين المعالج
    والذاكرة ووحدات الإدخال والإخراج. كل سلك ينقل بتًا واحدًا في كل مرة.
\end{tarif}

\begin{center}
\begin{tikzpicture}
    \node[unit, minimum width=2.8cm, minimum height=4cm] (cpu) at (5,0) {المعالج};
    \node[unit=brandPurple, minimum width=2.6cm, minimum height=1.2cm] (mem) at (0,2.6) {الذاكرة};
    \node[unit=brandGreen, minimum width=2.6cm, minimum height=1.2cm] (io) at (-4,2.6) {الإدخال\\والإخراج};
    % buses
    \draw[brandOrange, line width=5pt] (3.6,0.9) -- (-6,0.9);
    \draw[brandTeal!80!black, line width=5pt] (3.6,0) -- (-6,0);
    \draw[brandMuted, line width=5pt] (3.6,-0.9) -- (-6,-0.9);
    \node[font=\small\bfseries, lawn=orangeDeep, anchor=east] at (-6.1,0.9) {حافلة العناوين};
    \node[font=\small\bfseries, lawn=tealDeep, anchor=east] at (-6.1,0) {حافلة المعطيات};
    \node[font=\small\bfseries, lawn=brandMuted, anchor=east] at (-6.1,-0.9) {حافلة التحكم};
    \foreach \x/\n in {0/mem, -4/io} {
        \draw[brandOrange, line width=2pt, -{Stealth}] (\x-0.6,0.9) -- (\x-0.6,1.95);
        \draw[brandTeal!80!black, line width=2pt, {Stealth}-{Stealth}] (\x,0) -- (\x,1.95);
        \draw[brandMuted, line width=2pt, {Stealth}-{Stealth}] (\x+0.6,-0.9) -- (\x+0.6,1.95);
    }
\end{tikzpicture}
\end{center}

\begin{center}
\begin{tabularx}{\linewidth}{|>{\centering\arraybackslash}p{3cm}|>{\centering\arraybackslash}p{2.6cm}|R|}
    \hline
    \ras{الحافلة} & \ras{الاتجاه} & \ras{ماذا تنقل؟} \\ \hline
    \textbf{العناوين} \textenglish{(Address bus)} & من المعالج فقط & رقم الخانة التي يريد المعالج قراءتها أو الكتابة فيها \\ \hline
    \textbf{المعطيات} \textenglish{(Data bus)} & في الاتجاهين & القيمة المقروءة أو المكتوبة \\ \hline
    \textbf{التحكم} \textenglish{(Control bus)} & في الاتجاهين & إشارات مثل: قراءة \textenglish{(RD)}، كتابة \textenglish{(WR)}، الساعة، المقاطعة \\ \hline
\end{tabularx}
\end{center}

\begin{tashbih}
    تخيّل أنك تطلب طردًا من مكتب البريد: \textbf{العنوان} يخبر الموظف أي رف يقصد (حافلة العناوين)،
    و\textbf{الطرد} نفسه هو المعطاة (حافلة المعطيات)، وكلمة «أعطني» أو «خذ» هي إشارة التحكم
    (قراءة أو كتابة).
\end{tashbih}

\begin{khatwa}[: المعالج يقرأ الخانة \textenglish{0003}]
    \begin{enumerate}
        \item يضع المعالج العنوان \textenglish{0003} على \textbf{حافلة العناوين}.
        \item يرسل إشارة \textbf{قراءة} \textenglish{(RD)} على \textbf{حافلة التحكم}.
        \item تضع الذاكرة محتوى الخانة (\hex{41}) على \textbf{حافلة المعطيات}.
        \item يأخذ المعالج القيمة \hex{41} ويحفظها في أحد سجلاته.
    \end{enumerate}
\end{khatwa}

\subsection{عرض الحافلات}
\begin{itemize}
    \item \textbf{عرض حافلة العناوين} يحدد \textbf{سعة الذاكرة القصوى}: $n$ خطًا تعطي $2^n$ عنوانًا.
    \item \textbf{عرض حافلة المعطيات} يحدد \textbf{كم بتًا ننقل في كل مرة}: حافلة من 16 بتًا
          تنقل بايتين في المرة الواحدة.
\end{itemize}

\begin{khatwa}[: نقل 1000 بايت بحافلات مختلفة]
    \begin{enumerate}
        \item حافلة معطيات من \textbf{8 بتات} تنقل بايتًا واحدًا في كل مرة: نحتاج إلى $1000$ عملية نقل.
        \item حافلة من \textbf{16 بتًا} تنقل بايتين: $1000 \div 2 = 500$ عملية.
        \item حافلة من \textbf{64 بتًا} تنقل 8 بايتات: $1000 \div 8 = 125$ عملية فقط.
    \end{enumerate}
    كلما كانت الحافلة أعرض، انتقلت المعطيات أسرع، كما يمر عدد أكبر من السيارات في طريق سيّار
    من 8 مسالك مقارنة بطريق من مسلك واحد.
\end{khatwa}

\begin{center}
\begin{tabular}{|c|c|c|c|}
    \hline
    \ras{المعالج} & \ras{حافلة المعطيات} & \ras{حافلة العناوين} & \ras{الذاكرة القصوى} \\ \hline
    \textenglish{Intel 8086} & 16 بتًا & 20 بتًا & \textenglish{$2^{20}$ = 1 MB} \\ \hline
    \textenglish{Intel 80386} & 32 بتًا & 32 بتًا & \textenglish{$2^{32}$ = 4 GB} \\ \hline
    \textenglish{x86-64} الحديث & 64 بتًا & 48 بتًا (عمليًا) & \textenglish{$2^{48}$ = 256 TB} \\ \hline
\end{tabular}
\end{center}

%================================================
\section{تمارين}
%================================================

\begin{tamrin}
    حوّل الأعداد الآتية: (أ) \hex{3C} إلى العشري والثنائي. (ب) العدد العشري 200 إلى الست عشري.
\end{tamrin}
\begin{hall}
    (أ) $\hex{3C} = 3 \times 16 + 12 = 60$، وبالثنائي $\bin{0011\,1100}$.\\
    (ب) $200 = 12 \times 16 + 8$، والرقم 12 يُكتب \textenglish{C}، إذن $200 = $ \hex{C8}.
\end{hall}

\begin{tamrin}
    معالج له حافلة عناوين من 24 بتًا، وكل خانة ذاكرة تحتوي على بايت. ما السعة القصوى للذاكرة؟
\end{tamrin}
\begin{hall}
    $2^{24} = 2^{4} \times 2^{20} = 16 \times 1\ \text{MB} = 16\ \text{MB}$.
\end{hall}

\begin{tamrin}
    نريد ذاكرة سعتها \textenglish{4 GB}. كم خطًا تحتاج حافلة العناوين؟
\end{tamrin}
\begin{hall}
    $4\ \text{GB} = 2^{2} \times 2^{30} = 2^{32}$ بايت، إذن نحتاج إلى \textbf{32 خطًا}.
\end{hall}

\begin{tamrin}
    أجب بصحيح أو خطأ مع التعليل:
    (أ) في نموذج فون نيومان، يُحفظ البرنامج في ذاكرة مختلفة عن المعطيات.
    (ب) حافلة العناوين تنقل المعلومات في الاتجاهين.
    (ج) ذاكرة \textenglish{DRAM} تحتفظ بمحتواها بعد إطفاء الحاسوب.
\end{tamrin}
\begin{hall}
    (أ) \textbf{خطأ}: في نموذج فون نيومان ذاكرة واحدة للبرنامج والمعطيات؛ الفصل خاصية نموذج هارفارد.\\
    (ب) \textbf{خطأ}: المعالج وحده يرسل العناوين، فهي في اتجاه واحد.\\
    (ج) \textbf{خطأ}: \textenglish{DRAM} ذاكرة متطايرة، تُمحى عند انقطاع التيار.
\end{hall}

\begin{tamrin}
    معالج تردده \textenglish{2 GHz}، وتحتاج كل تعليمة في المتوسط إلى 4 دورات. كم تعليمة ينفذ في الثانية؟
\end{tamrin}
\begin{hall}
    عدد الدورات في الثانية $= 2 \times 10^{9}$. عدد التعليمات $= \dfrac{2 \times 10^{9}}{4} = 0.5 \times 10^{9}$،
    أي \textbf{500 مليون تعليمة في الثانية}.
\end{hall}
